\clearpage\SourceLeaf{400}\ArithmeticHead{393}{R}{GRUNDBEGRIFFE DER LEHRE VON DEN REIHEN.}
\begingroup
\makeatletter\tagsleft@true\let\veqno\@@leqno\makeatother
\begin{center}4.\end{center}
\noindent\hspace*{\parindent}Wenn die Reihe \(a',\,a'',\,a'''\) u. s. f. die obere Grenze \(L'\) hat, so sieht man leicht, dass dieselbe auch mit Weglassung des ersten Gliedes, oder die Reihe \(a'',\,a'''\) u. s. f. eine obere Grenze \(L''\) haben müsse, und zwar wird \(L''\) kleiner als \(L'\) sein, wenn \(a'=L'\) und die folgenden Glieder alle wenigstens um eine bestimmte Grösse kleiner sind; sonst wird sein \(L''=L'\), in keinem Falle aber \(L''>L'\). Auf gleiche Art wird die Reihe \(a''',\,a^{\mathrm{IV}}\) u. s. f. die obere Grenze \(L'''\) haben, die Reihe \(a^{\mathrm{IV}},\,a^{\mathrm{V}}\) u. s. f. die obere Grenze \(L^{\mathrm{IV}}\) u. s. f., und so werden alle diese obern Grenzen eine neue Reihe bilden \(L',\,L'',\,L''',\,L^{\mathrm{IV}}\) u. s. f., in welcher kein Glied grösser[\textsuperscript{*})] als das vorhergehende sein kann. So leitet man aus der ursprünglichen Reihe \(-1,\,-2,\,-3\) u. s. f. die Reihe \(-1,\,-2,\,-3\) u. s. f. als obere Grenzen Reihe ab, die mit jener selbst überein kommt; aus der Reihe
\[
\frac12,\ \frac23,\ \frac34\ \text{u. s. f.}
\]
\noindent diese \(1,\,1,\,1\) u. s. f. wo alle Glieder 1 sind; aus der Reihe
\begin{equation}
1,\ \frac12,\ -\frac14,\ -\frac38,\ -\frac1{16},\ \frac5{32},\ \frac7{64},\ -\frac3{128}\ \text{u. s. f.,}\tag*{[1]}
\end{equation}
\noindent welche aus Entwickelung des Bruches
\begin{equation}
\frac1{1-\frac12x+\frac12xx}\tag*{[2]}
\end{equation}
\noindent entsteht\textsuperscript{**}), diese
\[
1,\ \frac12,\ \frac5{32},\ \frac5{32},\ \frac5{32},\ \frac5{32},\ \frac7{64},\ \frac{23}{1024}\ \text{u. s. f.}
\]
\noindent Auf gleiche Art werden, wenn die Reihe \(a',\,a'',\,a''',\,a^{\mathrm{IV}}\) u. s. f. eine untere Grenze \(M'\) hat, auch die Reihen \(a'',\,a'''\) u. s. f., \(a''',\,a^{\mathrm{IV}}\) u. s. f., \(a^{\mathrm{IV}},\,a^{\mathrm{V}}\) u. s. f. u. s. f. untere Grenzen \(M'',\,M''',\,M^{\mathrm{IV}}\) u. s. f. haben, und diese werden eine neue Reihe \(M',\,M'',\,M''',\,M^{\mathrm{IV}}\) u. s. f. bilden, in welcher kein Glied kleiner sein kann als das vorhergehende.
\par\medskip\noindent\rule{.24\linewidth}{.4pt}\par\smallskip
{\small
\noindent\hspace*{\parindent}[\textsuperscript{*}) Die Handschrift hat: \textls[120]{kleiner}.]
\par\noindent\hspace*{\parindent}\textsuperscript{**}) Die Gründe, worauf die Formation dieser Reihen beruhet, werden anders wo vorgetragen werden.
\par}
\par\smallskip\noindent X1.\hfill50
\clearpage\endgroup

